\subsection{Сила тяжести}
%%%%%%%%%%%%%%%%%%%%%%%%%%%%%%%%%%%%%%%%%%%%%%%%%%%%
%%%%%%%%%%%%%%%%%%%%%%%%%%%%%%%%%%%%%%%%%%%%%%%%%%%%
\z{На каком расстоянии от поверхности Земли сила притяжения космического корабля к Земле в 100 раз меньше, чем на ее поверхности?}
%
{$h = 9R_\text{З}$}

%%%%%%%%%%%%%%%%%%%%%%%%%%%%%%%%%%%%%%%%%%%%%%%%%%%%
\z{Найдите ускорение свободного падения на поверхности Марса, если известно, что Марс примерно в 10 раз легче Земли, а радиус Марса примерно в 2 раза меньше радиуса Земли.}
%
{$g_\text{М}=\frac{2}{5}g$}

%%%%%%%%%%%%%%%%%%%%%%%%%%%%%%%%%%%%%%%%%%%%%%%%%%%%
\z{Два одинаковых однородных шара, соприкасаясь, притягивают друг друга с силой $F$.  Как изменится сила, если увеличить массу каждого шара в $n$ раз?  Плотность материала шаров не изменяется.}
%
{$n^{\frac{4}{3}}$}

%%%%%%%%%%%%%%%%%%%%%%%%%%%%%%%%%%%%%%%%%%%%%%%%%%%%
\z{Космическая ракета при старте с поверхности Земли движется вертикально вверх с ускорением $20\, \text{м/с}^2$. В кабине ракеты находится космонавт массой $80\, \text{кг}$. Чему в момент взлета равна сила тяжести, действующая на космонавта? А чему при этом равен его вес?} 
%
{$mg=800\,\text{Н};~P=2000\,\text{Н}$}

%%%%%%%%%%%%%%%%%%%%%%%%%%%%%%%%%%%%%%%%%%%%%%%%%%%%
\z[{\begin{wrapfigure}[5]{r}{0.3\textwidth}%
		\vspace{-10pt}%
		\begin{flushright}%
 			\includegraphics[width=0.3\textwidth]{3.2.5}%
		\end{flushright}%
	\end{wrapfigure}%
	}]%
{Сила гравитационного притяжения между двумя бильярдными шарами равна $F_0$. В точку $A$, находящуюся на одной линии с первыми двумя шарами и на равном расстоянии от них, помещают третий бильярдный шар. Найдите гравитационную силу, которая теперь будет действовать на каждый из шаров со стороны двух остальных.}
%
{$F_1=5F_0;$~$F_1=0;$~$F_3=5F_0$}

%%%%%%%%%%%%%%%%%%%%%%%%%%%%%%%%%%%%%%%%%%%%%%%%%%%%
\z[{\begin{wrapfigure}[5]{r}{0.3\textwidth}%
		\vspace{-15pt}%
		\begin{flushright}%
 			\includegraphics[width=0.3\textwidth]{3.2.6}%
		\end{flushright}%
	\end{wrapfigure}%
	}]%
{Четырьмя натянутыми нитями груз закреплён на тележке. Сила натяжения горизонтальных нитей соответственно $T_1$ и $T_2$, а вертикальных — $T_3$ и $T_4$. С каким ускорением тележка движется по горизонтальной плоскости?
}
%
{$a = \frac{T_2 - T_1}{T_4 - T_3}g$}

%%%%%%%%%%%%%%%%%%%%%%%%%%%%%%%%%%%%%%%%%%%%%%%%%%%%
\z[{\begin{wrapfigure}[4]{r}{0.3\textwidth}%
		\vspace{-10pt}%
		\begin{flushright}%
 			\includegraphics[width=0.3\textwidth]{3.2.7}%
		\end{flushright}%
	\end{wrapfigure}%
	}]%
{Верёвка выдерживает груз максимальной массы $m_1$ при его движении с некоторым ускорением, направленным вверх, и груз максимальной массы $m_2$ при его движении с ускорением, в $n$ раз меньшим первого по величине и направленным вниз (см. рисунок). Груз какой максимальной массы можно подвесить к верёвке в покое?
}
%
{$m_3 = \frac{(n + 1)(m_1 m_2)}{n m_1 + m_2}$}

